\section{Research-Track particle and clock ansatz}
\label{app:research-track}

\[
\boxed{
\text{SUBSTRATE INTERPRETATION / NON-ORTHODOX RESEARCH TRACK}
}
\]

\paragraph{Status.}
This section introduces candidate SST closures that are not consequences of
orthodox Euler hydrodynamics, Gross--Pitaevskii theory, or established particle
physics.  The purpose is to formulate a constrained and falsifiable research
programme while preserving the distinction between geometric definitions,
canonically calibrated SST quantities, and unverified particle assignments.

Suppose, as an additional hypothesis, that microscopic matter structures are
finite-core vortex knots of an incompressible underlying substrate.  A
predictive implementation must derive its equilibrium states from a declared
continuum energy and finite-thickness constraint.  Arbitrary pairwise
short-range potentials introduced only to prevent strand crossings are
excluded.

\subsection{Resolved finite-core energy functional}

Let
\begin{equation}
\gamma_K:
S^1\rightarrow\mathbb{R}^3,
\qquad
s\mapsto\bm{X}_K(s)
\end{equation}
be an arclength-parametrized centreline of knot type \(K\).  Its resolved tube
radius is defined geometrically by
\begin{equation}
a_{\rm core}(K)
=
\operatorname{reach}(\gamma_K)
=
\min\left[
\inf_s\frac{1}{\kappa(s)},
\frac{1}{2}d_{\rm dcsd}(\gamma_K)
\right],
\label{eq:RT-reach}
\end{equation}
where \(d_{\rm dcsd}\) is the minimum doubly-critical self-distance
\cite{CantarellaKusnerSullivan2002}.

A candidate resolved-tube energy is
\begin{align}
\mathcal{E}_{\rm eff}[\gamma_K]
={}&
E_{\rm SI}^{(a)}[\gamma_K,\Gamma_K]
+
T_{\rm eff}L[\gamma_K]
+
B_{\rm eff}\oint_{\gamma_K}\kappa^2\,\dd s
+
E_{\rm twist}[\gamma_K],
\label{eq:RT-resolved-functional}
\\
L[\gamma_K]
={}&
\oint_{\gamma_K}\dd s .
\end{align}

Here \(E_{\rm SI}^{(a)}\) is a finite-core regularization of the kinetic
self-induction energy.  It may be represented generally as
\begin{equation}
E_{\rm SI}^{(a)}
=
\frac{\rhoF\Gamma_K^2}{8\pi}
\oint_{\gamma_K}
\oint_{\gamma_K}
\bm{t}(s)\cdot\bm{t}(s')\,
\mathcal{K}_{a_{\rm core}}
\left(
\left|
\bm{X}_K(s)-\bm{X}_K(s')
\right|
\right)
\dd s\,\dd s',
\label{eq:RT-regularized-self-induction}
\end{equation}
where
\begin{equation}
\mathcal{K}_{a}(r)
\longrightarrow
\frac{1}{r}
\qquad
(r\gg a)
\end{equation}
and \(\mathcal{K}_{a}(0)\) is finite.  The regularization kernel must be
derived from, or fixed by, a declared core profile before mass predictions
are evaluated.

The coefficients have dimensions
\begin{equation}
[T_{\rm eff}]
=
\mathrm{J\,m^{-1}},
\qquad
[B_{\rm eff}]
=
\mathrm{J\,m}.
\end{equation}
A finite-core scaling ansatz may take the form
\begin{equation}
B_{\rm eff}
=
c_B T_{\rm eff}a_{\rm core}^2,
\label{eq:RT-bending-scaling}
\end{equation}
where \(c_B\) is dimensionless.  Dimensional analysis does not determine
\(c_B=1\); it must follow from a resolved cross-sectional calculation.

The minimization problem is therefore
\begin{equation}
\min_{\gamma_K\in K}
\mathcal{E}_{\rm eff}[\gamma_K]
\qquad
\text{subject to}
\qquad
\operatorname{reach}(\gamma_K)
\geq
a_{\rm core}.
\label{eq:RT-constrained-minimization}
\end{equation}

For a polygonal discretization \(V\), this constraint becomes
\begin{align}
\operatorname{MinRad}(v_i)
&\geq
a_{\rm core},
\\
\frac{1}{2}d(p,q)
&\geq
a_{\rm core}.
\end{align}
The active curvature constraints are kinks and the active self-distance
constraints are struts.  Their non-negative KKT multipliers supply a
contact-stress skeleton for the resolved tube
\cite{AshtonCantarellaPiatekRawdon2011}.

The hard reach constraint is a geometric admissibility condition, not an
auxiliary Lennard--Jones interaction.  Bending stiffness can oppose
homothetic collapse, but it cannot by itself exclude a finite-curvature
self-crossing.

\subsection{Quantized circulation and characteristic velocity}

The circulation is retained in the canonical SST form
\begin{equation}
\Gamma_K
=
n_K\Gamma_0,
\qquad
n_K\in\mathbb{Z},
\label{eq:RT-circulation-quantization}
\end{equation}
with
\begin{equation}
\Gamma_0
=
2\pi\rc\vchar
=
9.68361920\times10^{-9}\,
\mathrm{m^2\,s^{-1}}.
\label{eq:RT-circulation-quantum}
\end{equation}

Thus the conserved circulation must not vary continuously with the knot
radius \(R_K\).  Instead, the characteristic resolved velocity may depend
on geometry through
\begin{equation}
U_K
=
\frac{\Gamma_K}{2\pi R_K}
\,
\mathcal{G}
\left(
\frac{a_{\rm core}}{R_K},
\mathcal{T}_K,
\mathfrak{q}_K
\right),
\label{eq:RT-characteristic-speed}
\end{equation}
where \(\mathcal{G}\) is dimensionless and \(\mathfrak{q}_K\) denotes any
independently specified internal coloring data.

Equation~\eqref{eq:RT-characteristic-speed} separates the canonical
circulation quantum from the geometry-dependent relation between circulation
and a chosen characteristic velocity.

\subsection{Topology-first particle assignment}

A predictive particle dictionary must be fixed before the observed mass
spectrum is used.  Let
\begin{equation}
\mathsf{A}(P)
=
\left(
K_P,
n_P,
\mathfrak{q}_P,
\text{framing data},
\text{component structure}
\right)
\label{eq:RT-assignment-map}
\end{equation}
denote the preregistered assignment for particle \(P\).

If a non-Abelian internal coloring is introduced, the mathematically
appropriate structure is a group homomorphism
\begin{equation}
\mathfrak{q}_K:
\pi_1
\left(
\mathbb{R}^3\setminus K
\right)
\longrightarrow
G,
\label{eq:RT-coloring-homomorphism}
\end{equation}
whose values on Wirtinger generators satisfy the knot-group crossing
relations.

A tuple such as
\begin{equation}
(L\bmod3,S\bmod2,\chi)
\end{equation}
may serve as a collection of discrete labels, but it is not by itself a
homomorphism into
\(\mathfrak{su}(3)\oplus\mathfrak{su}(2)\oplus\mathfrak{u}(1)\).

The trefoil \(K_{2,3}=3_1\) may be retained as the preregistered electron
benchmark of the current SST taxonomy.  Baryon assignments remain under
audit.  A torus-knot candidate \(K_{3,q}\), a three-component colored
composite, or another explicitly specified topology may be tested, but
threefold winding must not be identified with QCD color solely because
both involve the integer three.

\subsection{Preregistered falsifiability test}

Before fitting any mass, the following objects must be fixed:

\begin{enumerate}
\item the particle-to-topology map \(\mathsf{A}\);
\item the core regularization kernel \(\mathcal{K}_{a}\);
\item the circulation numbers \(n_P\);
\item the dimensionless coefficient \(c_B\);
\item the twist-energy law, if retained;
\item the admissible embedding class and numerical convergence criteria.
\end{enumerate}

Only one overall energy normalization may then be calibrated using the
electron mass.  Define the dimensionless minimized energy
\begin{equation}
\widehat{\mathcal{E}}_P
=
\min_{\gamma\in\mathcal{C}_P}
\widehat{\mathcal{E}}_{\rm eff}
\left[
\gamma;
\mathsf{A}(P)
\right],
\label{eq:RT-dimensionless-minimum}
\end{equation}
where \(\mathcal{C}_P\) is the preregistered admissible class.

The out-of-sample proton prediction is
\begin{equation}
\widehat{R}_{p/e}
=
\frac{
\widehat{\mathcal{E}}_p
}{
\widehat{\mathcal{E}}_e
}.
\label{eq:RT-predicted-proton-ratio}
\end{equation}
The model must reproduce
\begin{equation}
\widehat{R}_{p/e}
\simeq
1836.15267
\label{eq:RT-proton-target}
\end{equation}
within a preregistered numerical and theoretical uncertainty, without
changing the topology assignment, the kernel, or any dimensionless
coefficient after inspecting the result \cite{Mohr2025CODATA}.

The same fixed model must then predict additional ratios, including the
neutron-to-electron ratio and at least one further particle sector.  Failure
of any preregistered primary target falsifies the corresponding assignment
or the energy closure.

\subsection{Clock-rate discipline}

Define
\begin{equation}
\beta_K
=
\frac{U_K}{c},
\qquad
\Pi_a
=
\frac{a_{\rm core}}{R_K},
\qquad
\Pi_\kappa
=
a_{\rm core}^2\langle\kappa^2\rangle.
\label{eq:RT-clock-parameters}
\end{equation}

The canon-compatible local clock factor is
\begin{equation}
\frac{\dd\tau_K}{\dd T}
=
\sqrt{1-\beta_K^2}.
\label{eq:RT-canonical-clock}
\end{equation}

A Research-Track deviation may only be introduced in the explicitly
factorized form
\begin{equation}
\frac{\dd\tau_K}{\dd T}
=
\sqrt{1-\beta_K^2}
\left[
1+
\varepsilon_{\rm RT}
\Delta_K
\left(
\Pi_a,
\Pi_\kappa,
\mathcal{T}_K,
\mathfrak{q}_K,
\ldots
\right)
\right],
\label{eq:RT-clock-deviation}
\end{equation}
subject to
\begin{equation}
\Delta_K(0,0,\ldots)=0.
\end{equation}

The correction must be local, universal, independently derived, and
preregistered before experimental comparison.  A ratio involving the
galactic averaging radius \(R_{\rm gal}\) is not admitted into the local
clock law unless a specific non-local coupling mechanism is derived.

In the absence of such a derivation,
\begin{equation}
\varepsilon_{\rm RT}=0
\end{equation}
and Eq.~\eqref{eq:RT-canonical-clock} remains the operative clock law.

\subsection{Reconnection, framing, and twist}

A smooth filled Euler vortex tube has frozen topology and does not possess
the Gross--Pitaevskii reconnection mechanism.  The model must therefore
choose among:

\begin{enumerate}
\item exact topological conservation in the ideal limit;
\item a declared finite-scale nonideal reconnection law;
\item a compressible or order-parameter core theory.
\end{enumerate}

A framed finite-core tube may carry twist.  Any twist term
\(E_{\rm twist}\) must specify the material frame, its transport equation,
and the partition of helicity between linking, writhe, and twist
\cite{Moffatt1969}.  Candidate falsifiers include Kelvin-wave dispersion,
reconnection scaling, and helicity-transfer measurements.

\section{Minimal numerical programme}
\label{app:minimal-numerical-programme}

A staged numerical test must separate geometry, hydrodynamics, topology
assignment, and clock interpretation.

\begin{enumerate}

\item
Preregister the candidate particle topologies, circulation integers, core
profile, regularization kernel, and all dimensionless coefficients.

\item
Generate centreline discretizations for the preregistered knot and link
classes.  Measure
\[
L,\qquad
\operatorname{MinRad},\qquad
d_{\rm dcsd},\qquad
\operatorname{reach},
\qquad
\int\kappa^2\,\dd s.
\]

\item
Minimize Eq.~\eqref{eq:RT-resolved-functional} subject to the hard
finite-thickness conditions
\[
\operatorname{MinRad}\geq a_{\rm core},
\qquad
\frac12d_{\rm dcsd}\geq a_{\rm core}.
\]
Record the active struts, kinks, KKT multipliers, and constrained-gradient
residual.

\item
Perform core-resolution and discretization convergence tests.  The predicted
energy ratios must be stable under changes in vertex count, quadrature order,
kernel resolution, and optimizer tolerance.

\item
Calibrate only the overall energy scale on the electron.  Evaluate all
preregistered baryon and additional particle ratios without further fitting.

\item
For a fully contained closed vortex loop, verify the exact global identity
\begin{equation}
\int_{V}\bm{\omega}\,\dd V
=
\Gamma_K\oint_{\gamma_K}\dd\bm{X}
=
\bm{0}.
\label{eq:RT-closed-loop-zero-monopole}
\end{equation}

\item
Test the local coarse-grained relation
\begin{equation}
\overline{\omega}_z(\bm{x})
\simeq
\Gamma_K
\mathcal{L}(\bm{x})
p_z(\bm{x})
\label{eq:RT-local-polarization}
\end{equation}
only as a local segment average.  For fully contained closed loops, verify
that positive axial-vorticity regions are balanced by return-vorticity
regions.

\item
Compare separately:
\begin{enumerate}
\item straight vortex lines crossing the averaging boundary;
\item boundary-anchored or periodic vortex lines;
\item fully closed vortex knots;
\item space-filling vortex-cell models.
\end{enumerate}
Only the first two classes can support a nonzero global vector-vorticity
monopole without an explicit return region.

\item
Measure finite-size and homogenization errors as
\[
\frac{L_{\rm cg}}{R_{\rm gal}}\rightarrow0,
\qquad
\frac{\ell_v}{L_{\rm cg}}\rightarrow0.
\]

\item
Use Gross--Pitaevskii simulations for depleted-core quantum superfluids and
Euler, vortex-blob, or resolved finite-core methods for filled classical
tubes.  Do not transfer reconnection laws between these models without a
declared bridge.

\item
Do not introduce a modified clock law until the circulation, velocity,
energy, and topology closures have passed their independent tests.

\end{enumerate}

